\section{Number theory}

\entry{Modular basics}{$O(\log)$ each}{%
\code{power(a, b, m)} is $a^b \bmod m$; it needs $m < 2^{31}$ or the \code{a*a} inside
overflows --- for bigger moduli cast to \code{i128}.\\[0.3ex]
\code{ext\_gcd(a, b, x, y)} returns $g=\gcd$ and fills $ax+by=g$. \code{inv(a, m)}
gives $a^{-1} \bmod m$ for \emph{any} \code{m} with $\gcd(a,m)=1$, and \code{-1} when
none exists. When \code{m} is prime, \code{power(a, m-2, m)} is the same thing and
shorter.\\[0.3ex]
$a^b \bmod m$ with a \emph{gigantic} $b$: read $b$ as a string, reduce it mod
$\varphi(m)$ --- but only add $\varphi(m)$ back (i.e. use $b \bmod \varphi(m) +
\varphi(m)$) when the true $b \ge \varphi(m)$, otherwise you get the wrong answer for
small $b$. That correction is needed exactly when $\gcd(a,m) \ne 1$.}

%LABEL modbasics
\begin{lstlisting}
ll power(ll a, ll b, ll mod) {         // a^b % mod
    a %= mod; ll res = 1;
    for (; b > 0; b >>= 1, a = a * a % mod)
        if (b & 1) res = res * a % mod;
    return res;
}
ll ext_gcd(ll a, ll b, ll &x, ll &y) {     // a*x + b*y = g
    if (b == 0) { x = 1; y = 0; return a; }
    ll x1, y1;
    ll g = ext_gcd(b, a % b, x1, y1);
    x = y1;
    y = x1 - (a / b) * y1;
    return g;
}
ll inv(ll a, ll m) {                 // -1 if none exists
    ll x, y;
    if (ext_gcd(a, m, x, y) != 1) return -1;
    return ((x % m) + m) % m;
}
\end{lstlisting}

\entry{Fast power}{Ruiming \quad $O(\log p)$}{%
\code{fast(x, p, m)} is $x^p \bmod m$. Needs $m < 2^{31}$ or \code{x*x} overflows.
\emph{Euler's theorem}, \emph{Binomials (inverse)} below call it.}

%LABEL fastmod
%FROM Ruiming/math/Fastmod.cpp
\begin{lstlisting}
ll fast(ll x,ll p,ll m){
    x%=m;
    ll ans=1;
    while(p>0){
        if(p&1){
            ans*=x;
            ans%=m;
        }
        x*=x;
        x%=m;
        p>>=1;
    }
    return ans;
}
\end{lstlisting}

\entry{Fast gcd}{Ruiming}{%
Binary (Stein) gcd with \code{\_\_builtin\_ctz}. Both arguments must be
\textbf{non-zero}: \code{\_\_builtin\_ctz(0)} is undefined.}

%FROM Ruiming/math/Fast_gcd.cpp
\begin{lstlisting}
ll gcd(ll a,ll b)
{
    ll az=__builtin_ctz(a),bz=__builtin_ctz(b);
    ll z=min(az,bz);
    ll dif;
    b>>=bz;
    while(a)
    {
        a>>=az;
        dif=b-a;
        az=__builtin_ctz(dif);
        if(a<b) b=a;
        if(dif<0) a=-dif;
        else a=dif;
    }
    return b<<z;
}
\end{lstlisting}

\entry{CRT}{Ruiming \quad $O(\sum m_i)$}{%
A whole program, \textbf{1-indexed}: reads \code{n}, then \code{n} pairs
\code{m[i] a[i]} meaning $x \equiv a_i \pmod{m_i}$, prints the smallest such $x$.
Moduli must be \textbf{pairwise coprime}. Merges one congruence at a time by trying
\code{j = 0..m[i]-1}, so it is linear in the moduli --- use \emph{CRT (general)} when
they are large or not coprime.}

%FROM Ruiming/math/CRT.cpp
\begin{lstlisting}
ll n,lst[MAXN],a[MAXN],m[MAXN],ans[MAXN],k,temp;

int main(){
    ios_base::sync_with_stdio(false);
    cin.tie(NULL);
    cin>>n;


    for(int i=1;i<=n;i++){
        cin>>m[i]>>a[i];
    }

    ans[1]=a[1];
    for(int i=2;i<=n;i++){
        for(ll j=0;j<=m[i]-1;j++){
            if(((j*m[i-1]+ans[i-1])%m[i])==a[i]){
                ans[i]=j*m[i-1]+ans[i-1];
                break;
            }
        }
        m[i]*=m[i-1];
    }
    cout<<ans[n]<<'\n';
}
\end{lstlisting}

\entry{CRT (general)}{Martin \quad $O(\log)$}{%
Merges $x \equiv r_1 \pmod{m_1}$ and $x \equiv r_2 \pmod{m_2}$ into one congruence
modulo $\mathrm{lcm}(m_1,m_2)$. Moduli need \emph{not} be coprime; returns
\code{\{-1,-1\}} when the two are incompatible. Fold a list by merging one at a
time.\\[0.3ex]
Watch the size: the intermediate \code{m1 / g * m2} must fit, so with several moduli
near $10^9$ switch to \code{i128}.}

%LABEL crt
%DEP modbasics
\begin{lstlisting}
// x = r1 (mod m1), x = r2 (mod m2)  ->  {r, lcm} or {-1,-1}
pair<ll,ll> crt(ll r1, ll m1, ll r2, ll m2) {
    ll x, y, g = ext_gcd(m1, m2, x, y);
    if ((r2 - r1) % g) return {-1, -1};
    ll l = m1 / g * m2, t = m2 / g;
    ll k = ((r2 - r1) / g % t) * (x % t) % t;
    ll r = (r1 + (i128)m1 * k) % l;
    return {(r % l + l) % l, l};
}
\end{lstlisting}

\entry{Sieves}{$O(n \log\log n)$ \quad mem $n$}{%
\code{sieve(n)} lists the primes up to \code{n} inclusive. \code{spf(n)} gives the
smallest prime factor of every number, which factorises any \code{x <= n} in
$O(\log x)$ by dividing repeatedly --- that is the one to build when you must factorise
many numbers.}

%LABEL sieves
\begin{lstlisting}
vector<int> sieve(int n) {
    vector<char> comp(n + 1, 0); vector<int> ps;
    for (int i = 2; i <= n; i++) if (!comp[i]) {
        ps.push_back(i);
        for (ll j = (ll)i * i; j <= n; j += i) comp[j] = 1;
    }
    return ps;
}
vector<int> spf(int n) {         // smallest prime factor
    vector<int> s(n + 1);
    iota(all(s), 0);
    for (int i = 2; (ll)i * i <= n; i++) if (s[i] == i)
        for (int j = i * i; j <= n; j += i)
            if (s[j] == j) s[j] = i;
    return s;
}
\end{lstlisting}

\entry{Euler sieve}{Ruiming \quad $O(n)$}{%
A whole program: reads \code{n q}, sieves the primes up to \code{n} into
\code{primes} (linear sieve), then answers \code{q} queries ``the $a$-th prime''.
\code{MAXN} must exceed \code{n}.}

%FROM Ruiming/math/Euler_sieve.cpp
\begin{lstlisting}
vector<ll> primes;
bool p[MAXN];
ll n,q;

int main(){
    ios_base::sync_with_stdio(false);
    cin.tie(NULL);

    cin>>n>>q;

    for(ll i=2;i<=n;i++){
        if(!p[i]){
            primes.push_back(i);
        }

        for(ll j=0;j<primes.size();j++){
            if(primes[j]*i>n) break;
            p[primes[j]*i]=1;
            if(i%primes[j]==0) break;
        }
    }

    for(int i=1;i<=q;i++){
        ll a;
        cin>>a;
        cout<<primes[a-1]<<'\n';
    }
}
\end{lstlisting}

\entry{Linear sieve + Euler phi}{Ruiming \quad $O(n)$ \quad mem $n$}{%
A whole program: reads \code{n}, fills \code{phi[1..n]} and \code{primes} (the
Euler / linear sieve, each composite crossed out exactly once). Drop the
\code{phi} lines for a plain linear sieve. \code{MAXN} must exceed \code{n}.\\[0.3ex]
$\varphi(n) = n \prod_{p \mid n} (1 - 1/p)$; $\sum_{d \mid n} \varphi(d) = n$. Number of
divisors and their sum are multiplicative, so the same sieve shape computes them.\\[0.3ex]
$a^b \bmod m$ with a gigantic $b$: see \emph{Euler's theorem} below.}

%FROM Ruiming/math/Linearphi.cpp
\begin{lstlisting}
ll n,phi[MAXN],ans=1;
bool p[MAXN];
vector<ll> primes;

int main(){
    ios_base::sync_with_stdio(false);
    cin.tie(NULL);
    cin>>n;

    if(n==1) cout<<0<<'\n';
    else{
        phi[1]=1;
        for(int i=2;i<=n;i++){//modiying euler sieve
            if(!p[i]){
                primes.push_back(i);
                phi[i]=i-1;
            }

            for(int j=0;j<primes.size();j++){
                if(primes[j]*i>n) break;
                p[primes[j]*i]=1;
                if(i%primes[j]==0){
                    phi[i*primes[j]]=primes[j]*phi[i];
                    break;
                }else{
                    phi[i*primes[j]]=(primes[j]-1)*phi[i];
                }
            }
        }
    }
}
\end{lstlisting}

\entry{Euler phi of one number}{Ruiming \quad $O(\sqrt n)$}{%
\code{phi(n)} by trial division. For every value up to $n$ at once, use the linear
sieve above.}

%FROM Ruiming/math/Phi.cpp
\begin{lstlisting}
ll phi(ll n) {//Euler's phi function
    ll ans = 1;
    for(ll i = 2; i * i <= n; ++ i) {
        if(n % i == 0) {
            n /= i;
            ans *= i - 1;
            while(n % i == 0) {
                n /= i;
                ans *= i;
            }
        }
    }
    if(n > 1)//optimisation: n>1 means one prime is left
        ans *= n - 1;
    return ans;
}
\end{lstlisting}

\entry{Prime factorization}{Ruiming \quad $O(\sqrt n)$}{%
\code{factor(n)} returns the prime factors with multiplicity, in increasing order, by
trial division. Fine up to $\sim10^{12}$; beyond that use Pollard rho.}

%FROM Ruiming/math/Prime_factorization.cpp
\begin{lstlisting}
vector<ll> factor(ll n) {
    vector<ll> ret;
    for (ll i = 2; i * i <= n; i++) {
        while (n % i == 0) {
            ret.push_back(i);
            n /= i;
        }
    }
    if (n > 1) { ret.push_back(n); }
    return ret;
}
\end{lstlisting}

\entry{Factorization table}{Ruiming \quad $O(\texttt{MAXN}\sqrt{\texttt{MAXN}})$}{%
A whole program: \code{fac[i]} is the factorization of every \code{i < MAXN} as a
\code{map} prime $\rightarrow$ exponent, then reads \code{a[1..n]}. His file uses
\code{MAXN = 36}.}

%FROM Ruiming/math/Factorization.cpp
\begin{lstlisting}
ll n,a[MAXN],ans;
map<ll,ll> fac[MAXN];

int main(){
    ios_base::sync_with_stdio(false);
    cin.tie(NULL);
    cin>>n;

    for(int i=1;i<MAXN;i++){
        ll temp=i;
        for(int j=2;j<=sqrt(temp+0.5);j++){
            while(temp%j==0){
                fac[i][j]++;
                temp/=j;
            }
        }
        if(temp>1){
            fac[i][temp]++;
        }
    }

    for(int i=1;i<=n;i++){
        cin>>a[i];
    }


}
\end{lstlisting}

\entry{Euler's theorem (huge exponent)}{Ruiming \quad $O(m + |b|)$}{%
A whole program: reads \code{a m}, then the exponent \code{b} as a string, prints
$a^b \bmod m$. Computes $\varphi(m)$, reduces \code{b} mod $\varphi(m)$ while reading,
and adds $\varphi(m)$ back only if the true $b \ge \varphi(m)$ (\code{flag}) --- which
is what makes it correct when $\gcd(a,m) \ne 1$. Uses \code{fast} from
\emph{Fast power}.}

%DEP fastmod
%FROM Ruiming/math/Eulerthm.cpp
\begin{lstlisting}
bool flag;

ll gcd(ll x,ll y){
    while(x!=y){
        if(x>y) swap(x,y);
        if(y%x==0) return x;
        y-=x*(y/x);
    }
    return x;
}



ll a,m,phi=1;
string b;

ll read(){
    string x;
    cin>>x;
    ll ans=0;

    for(int i=0;i<x.size();i++){
        ans=(ans*10+(x[i]-'0'));
        if(ans>=phi) flag=true;
        ans%=phi;
    }
    return ans;
}

int main(){
    cin>>a>>m;
    //find the euler function to m

    ll d=gcd(a,m);

    ll temp=m;
    for(int i=2;i<=m;i++){
        //here uses temp still makes sure i is a prime
        if(temp%i==0){
            ll times=0;
            while(temp%i==0){
                times++;
                temp/=i;//makes sure i is a prime
            }
            phi*=pow(i,times-1)*(i-1);
        }
    }
    //if(phi==1) phi=m-1;

    //b needs to be stored as a string and takes mod phi
    //then we can calculate it with fast mod
    ll b=read();
    ll ans;
    if(flag) ans=fast(a,b+phi,m);
    else ans=fast(a,b,m);

    //deal with d d|m, will lead up to a ,
    //should be some period

    cout<<ans<<'\n';
}
\end{lstlisting}

\entry{Miller-Rabin (standalone)}{Ruiming \quad $O(12 \log n)$}{%
A complete file on its own --- \textbf{do not paste it under the header}: its
\code{typedef \_\_int128 ll} replaces \code{ll}, so \code{n} can exceed
\code{long long}, with \code{read}/\code{write}/\code{print} fast I/O for
\code{\_\_int128}. Reads \code{t} numbers, prints \code{Yes}/\code{No}. \code{a * a}
in \code{qpow} overflows once \code{n} passes $\sim1.3\cdot10^{19}$.}

%STANDALONE
%FROM Ruiming/math/miller_rabin.cpp
\begin{lstlisting}
#include <bits/stdc++.h>
using namespace std;

typedef __int128 ll;
typedef pair<int, int> pii;

template<typename T> inline T read() {
    T x = 0, f = 1; char ch = 0;
    for(; !isdigit(ch); ch = getchar()) if(ch == '-') f = -1;
    for(; isdigit(ch); ch = getchar())
        x = (x << 3) + (x << 1) + (ch - '0');
    return x * f;
}

template<typename T> inline void write(T x) {
    if(x < 0) putchar('-'), x = -x;
    if(x > 9) write(x / 10);
    putchar(x % 10 + '0');
}

template<typename T> inline void print(T x, char ed = '\n') {
    write(x), putchar(ed);
}

ll t, n;

ll qpow(ll a, ll b, ll mod) {
    ll ret = 1;
    while(b) {
        if(b & 1) ret = (ret * a) % mod;
        a = (a * a) % mod;
        b >>= 1;
    }
    return ret % mod;
}

vector<ll> p = {2, 3, 5, 7, 11, 13, 17, 19, 23, 29, 31, 37};

bool miller_rabin(ll n) {
    if(n < 3 || n % 2 == 0) return n == 2;
    ll u = n - 1, t = 0;
    while(u % 2 == 0) u /= 2, ++ t;
    for(auto a : p) {
        if(n == a) return 1;
        if(n % a == 0) return 0;
        ll v = qpow(a, u, n);
        if(v == 1) continue;
        ll s = 1;
        for(; s <= t; ++ s) {
            if(v == n - 1) break;
            v = v * v % n;
        }
        if(s > t) return 0;
    }
    return 1;
}

int main() {
    t = read<ll>();
    while(t --) {
        n = read<ll>();
        if(miller_rabin(n)) puts("Yes");
        else puts("No");
    }
    return 0;
}
\end{lstlisting}

\entry{Miller-Rabin + Pollard rho}{Ruiming \quad $O(n^{1/4})$}{%
\code{miller\_rabin(n)} is \textbf{deterministic} for every \code{n} that fits in
\code{long long} (those 12 bases suffice below $3.3\cdot10^{24}$).\\[0.3ex]
\code{p.clear(); pollard\_rho(n);} leaves the prime factors of \code{n} in the global
\code{p}, with multiplicity, unsorted --- sort to group them. Multiplications go through
\code{i128}, so \code{n} up to $9\cdot10^{18}$ is safe.\\[0.3ex]
For $n \le 10^{12}$ or so, plain trial division up to $\sqrt n$ is shorter.}

%LABEL pollard
%FROM Ruiming/math/Pollard_rho.cpp
\begin{lstlisting}
#define LL long long
#define i128 __int128
mt19937_64 rd(time(0));

i128 ksm(i128 x, i128 y, LL mod)
{
    i128 res = 1;
    x %= mod;
    while(y)
    {
        if(y & 1) res = res * x % mod;
        y >>= 1;
        x = x * x % mod;
    }
    return res;
}

bool miller_rabin(LL n)
{
    if(n <= 1 || n > 2 && ~n & 1) return false;
    LL m = n - 1, k = 0;
    while(~m & 1) m >>= 1, k++;
    vector<int> s = {2, 3, 5, 7, 11, 13, 17, 19, 23, 29,
                     31, 37};
    for(int a : s)
    {
        if(a >= n) break;
        i128 x = ksm(a, m, n), y;
        for(int i = 1; i <= k; i++, x = y)
        {
            y = x * x % n;
            if(y == 1 && x != 1 && x != n - 1) return false;
        }
        if(x != 1) return false;
    }
    return true;
}

vector<LL> p;
void pollard_rho(LL n)
{
    if(n <= 1) return;
    if(miller_rabin(n))
    {
        p.push_back(n);
        return;
    }
    LL d = rd() % (n - 1) + 1;
    auto f = [&](LL x)
    {
        return ((i128)x * x + d) % n;
    };
    LL x = rd() % (n - 1) + 1, y = f(x), t = 1;
    while(t == 1)
    {
        for(int i = 1; i <= 32; i++)
        {
            while(x == y)
                d = rd() % (n - 1) + 1,
                x = rd() % (n - 1) + 1, y = f(x);
            if((i128)t * abs(x - y) % n == 0) break;
            t = (i128)t * abs(x - y) % n;
            x = f(x), y = f(f(y));
        }
        t = __gcd(n, t);
    }
    pollard_rho(t), pollard_rho(n / t);
}
\end{lstlisting}

\entry{Binomials mod p}{precompute $O(n)$, query $O(1)$ \quad mem $n$}{%
Call \code{precompute()} once, then \code{C(n, k)}. Uses the header's \code{MAXN}, so
\textbf{raise \code{MAXN}} above the largest \code{n}. \code{MOD} must be
\textbf{prime} and larger than \code{MAXN}. Returns 0 outside the triangle, so no
guards at the call site.\\[0.3ex]
Non-prime modulus, or $n$ small: build Pascal's triangle instead
(\code{c[i][j] = c[i-1][j-1] + c[i-1][j]}), which needs no inverses. Huge $n$ with small
prime $p$: Lucas' theorem, multiply \code{C} of the base-$p$ digits.\\[0.3ex]
Stars and bars: ways to write $n$ as an ordered sum of $k$ non-negative parts is
\code{C(n+k-1, k-1)}. Catalan $C_n = $ \code{C(2n,n) * inv(n+1)}.}

%DEP modbasics
%LABEL binom
\begin{lstlisting}
const ll MOD = 1000000007;      // MAXN from the header
ll fact[MAXN], inv_fact[MAXN];

void precompute() {
    fact[0] = 1;
    rep(i, 1, MAXN) fact[i] = fact[i-1] * i % MOD;
    inv_fact[MAXN-1] = power(fact[MAXN-1], MOD - 2, MOD);
    down(i, MAXN - 1, 0)
        inv_fact[i] = inv_fact[i+1] * (i+1) % MOD;
}
ll C(ll n, ll k) {
    if (n < 0 || k < 0 || k > n || n >= MAXN) return 0;
    return fact[n] * inv_fact[k] % MOD
                   * inv_fact[n-k] % MOD;
}
\end{lstlisting}

\entry{Binomials (Pascal)}{Ruiming \quad $O(\texttt{MAXN} \cdot 106)$}{%
A whole program: Pascal's triangle \code{c[i][j]} mod \code{M} for \code{i < MAXN-1},
\code{j < 106}. Needs no inverses, so it works for any modulus.}

%FROM Ruiming/math/Combinumber.cpp
\begin{lstlisting}
const ll M=1000000007;
ll c[MAXN][106],t,n,m,k,lst[MAXN],ans;

int main(){
    ios_base::sync_with_stdio(false);
    cin.tie(NULL);

    c[0][0]=1;

    for(int i=1;i<MAXN-1;i++){
        c[i][0]=1;
        if(i<106) c[i][i]=1;
        for(int j=1;j<min(i,105);j++){
            c[i][j]=c[i-1][j]+c[i-1][j-1];
            c[i][j]%=M;
        }
    }
}
\end{lstlisting}

\entry{Binomials (inverse)}{Ruiming \quad $O(n \log \texttt{MOD})$}{%
A whole program: reads \code{n k}, fills \code{fac[i]} $= i!$ and \code{inv[i]}
$= (i!)^{-1}$ mod \code{MOD}, then $\binom{n}{i}$ is \code{comb} in the last loop. Uses
\code{fast} from \emph{Fast power}. His file uses \code{MAXN = 26}.}

%DEP fastmod
%FROM Ruiming/math/Combinumber_inv.cpp
\begin{lstlisting}
const ll MOD=1e9+7;
ll t,n,fac[MAXN],inv[MAXN],ans,k;

int main(){
    ios_base::sync_with_stdio(false);
    cin.tie(NULL);
    cin>>n>>k;
    ll temp=1;
    for(int i=1;i<=n;i++){
        temp*=i;
        temp%=MOD;
        inv[i]=fast(temp,MOD-2,MOD);
        fac[i]=temp;
    }

    fac[0]=1;
    inv[0]=1;

    for(int i=0;i<=n;i++){
        ll comb=((fac[n]*inv[i])%MOD*inv[n-i])%MOD;
    }
}
\end{lstlisting}

\entry{Linear inverses}{Ruiming \quad $O(n)$}{%
A whole program: reads \code{n p} ($p$ prime), prints $i^{-1} \bmod p$ for
$i = 1..n$, using one modular exponentiation and the identity
$i^{-1} = (i-1)! \cdot (i!)^{-1}$.}

%FROM Ruiming/math/Linearinverse.cpp
\begin{lstlisting}
ll n,p,ans[MAXN],corr[MAXN],now,fac[MAXN],inv[MAXN];

ll fastMod(ll a,ll n,ll mod){
    ll ans=1;
    a%=mod;
    while(n){
        if(n&1){
            ans=(ans*a)%mod;
        }
        a=(a*a)%mod;
        n>>=1;
    }
    return ans;
}

int main(){
    cin>>n>>p;

    fac[0]=1;
    for(int i=1;i<=n;i++){
        fac[i]=(fac[i-1]*i)%p;
    }

    inv[n]=fastMod(fac[n],p-2,p);//inv fac[i]
    for(int i=n-1;i>=1;i--){
        inv[i]=((i+1)*inv[i+1])%p;
    }

    for(int i=1;i<=n;i++){
        cout<<(fac[i-1]*inv[i])%p<<'\n';
    }
}
\end{lstlisting}

\entry{Matrix power}{Ruiming \quad $O(k^3 \log n)$}{%
\textbf{1-indexed}, fixed \code{m[MAXN][MAXN]}: set \code{row}/\code{col} and fill
\code{m[1..row][1..col]}. \code{a * b} needs \code{a.col == b.row};
\code{pow\_matrix(a, n)} needs \code{a} square.\\[0.3ex]
\textbf{Delete the header's \code{typedef vector<vi> matrix;}} --- same name. Keep
\code{MAXN} small (his file uses 106): every \code{matrix} is $8\,\texttt{MAXN}^2$ bytes
and \code{operator*} builds one on the stack. Clashes with \code{MOD} in
\emph{Binomials} if both are pasted.\\[0.3ex]
Linear recurrences: Fibonacci is \code{\{\{1,1\},\{1,0\}\}} to the $n$, entry
\code{m[1][2]}. With \code{min}/\code{+} instead of \code{+}/\code{*} it gives
``cheapest walk using exactly $n$ edges''.}

%UNDEF matrix
%FROM Ruiming/math/Matrix.cpp
\begin{lstlisting}
const ll MOD=1e9+7;

struct matrix{
    ll row,col;
    ll m[MAXN][MAXN];
};

matrix operator * (const matrix &a,const matrix &b){
    //assuming they fits a.col==b.row

    matrix c;
    memset(c.m,0,sizeof(c.m));
    c.row=a.row;
    c.col=b.col;
    for(int i=1;i<=c.row;i++){
        for(int j=1;j<=c.col;j++){
            for(int k=1;k<=a.col;k++){
                c.m[i][j]+=(a.m[i][k]*b.m[k][j])%MOD;
                c.m[i][j]%=MOD;
            }
        }
    }
    return c;
}

matrix pow_matrix(matrix a,ll n){
    //assuming a is a square matrix
    matrix res;
    memset(res.m,0,sizeof(res.m));
    res.row=res.col=a.col;
    for(int i=1;i<MAXN;i++) res.m[i][i]=1;

    while(n){
        if(n&1) res=res*a;
        a=a*a;
        n/=2;
    }

    return res;
}
\end{lstlisting}

\entry{Gaussian elimination}{Ruiming \quad $O(n^3)$}{%
A whole program, \textbf{1-indexed}, square $n\times n$: reads \code{n}, then each row
of \code{A} followed by its \code{b[i]}. Prints \code{-1} if there is no solution,
\code{0} if there are infinitely many, else \code{x1=..} to 2 decimals. Forward
elimination, then back substitution.\\[0.3ex]
Zero tests are exact (\code{==0}) and there is no largest-pivot choice: fine for integer
input, but with real input compare \code{fabs(..) < eps} instead.\\[0.3ex]
Over $\mathbb{F}_2$ use \code{bitset<N>} rows and \code{\textasciicircum=}, which is
64x faster and exact --- that is the version for XOR-basis and lights-out problems.
Mod a prime: replace \code{/} with \code{inv}.}

%FROM Ruiming/math/gaussjordanelimination.cpp
\begin{lstlisting}
ll n;
double A[MAXN][MAXN],b[MAXN],ans[MAXN];
stack<pair<ll,ll>> swaped;

int main(){
    ios_base::sync_with_stdio(false);
    cin.tie(NULL);
    cin>>n;
    for(int i=1;i<=n;i++){
        for(int j=1;j<=n;j++){
            cin>>A[i][j];
        }
        cin>>b[i];
    }

    ll col=1;
    for(int i=1;i<=n&&col<=n;i++,col++){//row
        while(A[i][col]==0&&col<=n){
            bool found=false;
            for(int j=i+1;j<=n;j++){
                if(A[j][col]!=0){
                    swap(A[j],A[i]);
                    swap(b[j],b[i]);//swap b as well
                    found=true;
                    break;
                }
            }
            if(!found) col++;
        }

        for(int j=i+1;j<=n;j++){
            if(A[i][col]==0) break;//excepet bug
            double fac=A[j][col]/A[i][col];
            for(int k=col;k<=n;k++){
                A[j][k]-=fac*A[i][k];
            }
            b[j]-=fac*b[i];
        }
    }

    bool no=false,p=false;
    for(int i=1;i<=n;i++){
        bool zero=true;

        for(int j=1;j<=n;j++){
            if(A[i][j]!=0){
                zero=false;
                break;
            }
        }
        p|=zero;

        if(zero&&b[i]!=0){
            no=true;
            break;
        }
    }

    if(p){
        if(no){
            cout<<-1<<'\n';
        }else cout<<0<<'\n';
    }else{
        for(int i=n;i>=1;i--){
            for(int j=i+1;j<=n;j++){
                b[i]-=ans[j]*A[i][j];
            }
            ans[i]=b[i]/=A[i][i];
        }
        for(int i=1;i<=n;i++){
            printf("x%d=%.2f\n",i,ans[i]);
        }
    }
}
\end{lstlisting}

\entry{Big integers}{Ruiming}{%
Decimal strings, optional leading \code{-}. \code{Bigadd}, \code{Bigsub},
\code{Bigmul} ($O(nm)$ schoolbook), \code{Bighalf} (non-negative only).
\code{Scompare}/\code{Ecompare} are $<$ / $\le$ for non-negative numbers.
\code{BigMod(x, m)} is $x \bmod m$ by binary search on the quotient (the same search
gives division); \code{BigModSmall(a, mod)} is a big number mod an \code{ll}.\\[0.3ex]
\code{Bigadd} of a positive and a negative number is \emph{not} handled (those two lines
are commented out) --- call \code{Bigsub} yourself.}

%FROM Ruiming/math/Big_Integer.cpp
\begin{lstlisting}
string Bigadd(string x,string y){
    //if(x[0]!='-'&&y[0]=='-')
    //    return Bigsub(x,y.substr(1,y.size()-1));
    //if(x[0]=='-'&&y[0]!='-')
    //    return Bigsub(y,x.substr(1,x.size()-1));
    if(x[0]=='-'&&y[0]=='-') return "-"+Bigadd(
        y.substr(1,y.size()-1),x.substr(1,x.size()-1));

    int s=max(x.size(),y.size()),a[s+6],b[s+6];
    int res[s+6];
    memset(res,0,sizeof(res));
    memset(a,0,sizeof(a));
    memset(b,0,sizeof(b));

    for(int i=0;i<x.size();i++) a[i]=x[x.size()-1-i]-'0';
    for(int i=0;i<y.size();i++) b[i]=y[y.size()-1-i]-'0';

    for(int i=0;i<s;i++){
        res[i]+=a[i]+b[i];
        res[i+1]=res[i]/10;
        res[i]%=10;
    }
    string ans="";
    bool first=false;
    for(int i=s+1;i>=0;i--){
        if(res[i]==0&&!first) continue;
        else{
            ans+=res[i]+'0';
            first=true;
        }
    }
    if(ans=="") return "0";
    return ans;
}

string Bigsub(string x,string y){
    if(x[0]!='-'&&y[0]=='-')
        return Bigadd(x,y.substr(1,y.size()-1));
    if(x[0]=='-'&&y[0]!='-')
        return '-' + Bigadd(x.substr(1,x.size()-1),y);
    if(x[0]=='-'&&y[0]=='-') return Bigsub(
        y.substr(1,y.size()-1),x.substr(1,x.size()-1));


    bool neg=false;
    if(y.size()>x.size()){
        swap(x,y);
        neg=true;
    }

    if(y.size()==x.size()&&y>x){
        swap(x,y);//ensures y>x
        neg=true;
    }

    int s=max(x.size(),y.size()),a[s+6],b[s+6];
    int res[s+6];
    memset(res,0,sizeof(res));
    memset(a,0,sizeof(a));
    memset(b,0,sizeof(b));

    for(int i=0;i<x.size();i++) a[i]=x[x.size()-1-i]-'0';
    for(int i=0;i<y.size();i++) b[i]=y[y.size()-1-i]-'0';

    for(int i=0;i<s;i++){
        res[i]+=a[i]-b[i];
        while(res[i]<0){
            res[i]+=10;
            res[i+1]--;
        }
    }
    string ans="";
    if(neg) ans+="-";
    bool first=false;
    for(int i=s+1;i>=0;i--){
        if(res[i]==0&&!first) continue;
        else{
            ans+=res[i]+'0';
            first=true;
        }
    }
    if(ans==""||ans=="-") return "0";
    return ans;
}

// multiplication not done like this
string Bigmul(string x,string y){
    if(x=="0"||y=="0") return "0";

    if(x[0]!='-'&&y[0]=='-')
        return "-"+Bigmul(x,y.substr(1,y.size()-1));
    if(x[0]=='-'&&y[0]!='-')
        return "-"+Bigmul(y,x.substr(1,x.size()-1));
    if(x[0]=='-'&&y[0]=='-'){
        x=x.substr(1,x.size()-1);
        y=y.substr(1,y.size()-1);
    }

    int s=max(x.size(),y.size()),a[2*s+6],b[2*s+6];
    int res[2*s+6];
    memset(res,0,sizeof(res));
    memset(a,0,sizeof(a));
    memset(b,0,sizeof(b));

    for(int i=0;i<x.size();i++) a[i]=x[x.size()-1-i]-'0';
    for(int i=0;i<y.size();i++) b[i]=y[y.size()-1-i]-'0';


    for(int i=0;i<x.size();i++){
        for(int j=0;j<y.size();j++){
            res[i+j]+=b[j]*a[i];
            res[i+j+1]+=res[i+j]/10;//will not increase 2.
            res[i+j]%=10;
        }
    }

    string ans="";
    bool first=false;
    for(int i=2*s+5;i>=0;i--){
        if(res[i]==0&&!first) continue;
        else{
            ans+=res[i]+'0';
            first=true;
        }
    }
    if(ans=="") return "0";
    return ans;
}

string Bighalf(string x){
    int s=x.size(),a[s+6];
    int res[s+6];
    memset(res,0,sizeof(res));
    memset(a,0,sizeof(a));

    for(int i=0;i<x.size();i++) a[i]=x[i]-'0';

    for(int i=0;i<x.size();i++){
        res[i]+=a[i]/2;
        if(a[i]%2){
            a[i+1]+=10;
        }
    }

    string ans="";
    bool first=false;
    for(int i=0;i<s;i++){
        if(res[i]==0&&!first) continue;
        else{
            ans+=res[i]+'0';
            first=true;
        }
    }
    if(ans=="") return "0";
    return ans;
}

bool Scompare(string x,string y){
    if(x.size()<y.size()) return true;
    if(x.size()>y.size()) return false;
    return x<y;
}

bool Ecompare(string x,string y){
    if(x.size()<y.size()) return true;
    if(x.size()>y.size()) return false;
    return x<=y;
}

//can get Bigdiv in the same way
string BigMod(string x,string m){
    string start="0",end="1";
    for(int i=1;i<=x.size()+6;i++) end+="0";

    while(Scompare(start,end)){
        string temp=Bigadd(end,"1");
        string mid=Bigadd(start,Bighalf(Bigsub(temp,start)));
        if(Ecompare(Bigmul(mid,m),x)){
            start=mid;
        }else{
            end=Bigsub(mid,"1");
        }
    }
    string res=Bigsub(x,Bigmul(start,m));
    return res;
}

ll BigModSmall(string a,ll mod){
    ll x=0;
    for(int i=0;i<a.size();i++){
        x=((10*x)+(a[i]-'0'))%mod;
    }
    return x;
}
\end{lstlisting}

\entry{FFT / NTT (convolution)}{$O(n \log n)$ \quad mem $n$}{%
\code{conv(a, b)} multiplies two polynomials: \code{c[i] = sum a[j]*b[i-j]}, length
\code{a.size()+b.size()-1}. This is also ``for every pair of elements, bucket their
sum'' --- count pairs summing to $s$, subset-sum counts, string matching with
wildcards.\\[0.3ex]
\textbf{NTT}, exact, mod \code{P = 998244353} (a prime with $2^{23} \mid P-1$, root 3).
So the result must be wanted mod that prime, and the transform length must be a power
of two $\le 2^{23}$. Other friendly primes: $1004535809$ (root 3),
$167772161$ (root 3).\\[0.3ex]
If the true coefficients are needed exactly and exceed \code{P}, run it under two
different primes and \code{crt} the results. For real/complex input use a
\code{complex<double>} FFT instead --- but it loses precision above $\sim10^{14}$ in the
output, so prefer NTT whenever the answer is modular.}

%LABEL ntt
\begin{lstlisting}
const ll P = 998244353, G = 3;
ll pw(ll a, ll b, ll m = P) {
    ll r = 1; a %= m;
    for (; b; b >>= 1, a = a * a % m)
        if (b & 1) r = r * a % m;
    return r;
}
void ntt(vi& a, bool inv) {
    int n = (int)a.size();
    for (int i = 1, j = 0; i < n; i++) {
        int b = n >> 1;
        for (; j & b; b >>= 1) j ^= b;
        j ^= b;
        if (i < j) swap(a[i], a[j]);
    }
    for (int len = 1; len < n; len <<= 1) {
        ll w = pw(G, (P - 1) / (2 * len));
        if (inv) w = pw(w, P - 2);
        for (int i = 0; i < n; i += 2 * len) {
            ll cur = 1;
            rep(j, 0, len) {
                ll u = a[i+j], v = a[i+j+len] * cur % P;
                a[i+j] = (u + v) % P;
                a[i+j+len] = (u - v + P) % P;
                cur = cur * w % P;
            }
        }
    }
    if (inv) {
        ll ni = pw(n, P - 2);
        for (ll& x : a) x = x * ni % P;
    }
}
vi conv(vi a, vi b) {
    if (a.empty() || b.empty()) return {};
    int need = (int)(a.size() + b.size() - 1), n = 1;
    while (n < need) n <<= 1;
    a.resize(n); b.resize(n);
    ntt(a, false); ntt(b, false);
    rep(i, 0, n) a[i] = a[i] * b[i] % P;
    ntt(a, true);
    a.resize(need);
    return a;
}
\end{lstlisting}
